\documentclass[a4paper, 11pt]{article}


\usepackage[slovak]{babel}
\usepackage[utf8]{inputenc}
\usepackage[left=2cm, top=3cm, text={17cm, 24cm}]{geometry}
\usepackage{times}
\usepackage{verbatim}
\usepackage{enumitem}
\usepackage{graphicx} 
\usepackage[unicode]{hyperref}
\usepackage{listings}
\usepackage{color}


\begin{document}

	\begin{titlepage}
		\begin{center}
			\Huge
			\textsc{Vysoké učení technické v~Brně} \\
			\huge
			\textsc{Fakulta informačních technologií} \\
			\vspace{\stretch{0.382}}
			\begin{figure}[h]
                \centering
                \scalebox{0.6}{\includegraphics{logo.eps}}
            \end{figure}
			\LARGE
			Projektová dokumentácia k~predmetu IMS \\
			\Huge
			Zadanie č. 3 - Epidemiologické modely - mikro-úroveň
			\vspace{\stretch{0.618}}
		\end{center}

		{\Large
			\today
			\hfill
			\begin{tabular}{r}
			Dávid Oravec (xorave05) \\
			Dominik Boboš (xbobos00)
			\end{tabular}
		}
	\end{titlepage}
	\newpage
	
	\tableofcontents
	\newpage
	
	\section{Úvod}
	\par Pre vytvorenie epidemiologického modelu, sme si vybrali ochorenie COVID-19. Je to infekčné ochorenie vyvolané koronavírusom SARS-CoV-2, ktorý patrí do rodu Betacoronavirus. Prvý oficiálne potrvdený prípad sa datuje na 8. december 2019, aj keď niektoré nepotvrdené zdroje uvádzajú 17. november 2019. \cite{guardian}
	\newline
	
	\par Ochorenie postihlo celý svet a tak po roku od prvého prípadu sa stále nachádza v pandemickej kríze. COVID-19 zasiahol do života ľudí a v boji proti tejto pandémii musia vlády všetkých krajín zavádzať opatrenia, ktoré pomôžu zabrániť šíreniu. Väčšina týchto opatrení má devastačné následky v oblasti ekonomiky, kultúry, športu, ale aj psychického zdravia ľudí.
	\newline
	
	\par V súvislosti s touto problematikou sme sa rozhodli zamerať na skúmanie šírenia koronavírusu v oblasti kultúry, konrétne kinematografie. Našim cieľom je zistiť, akým spôsobom sa vírus môže šíriť medzi ľuďmi, ktorí sú uzavretí v realtívne malom priestore (kinosále) a ako znížiť riziko nákazy bez potreby zavádzania drastických opatrení.
	
	\subsection{Zdroje faktov}
	    \subsubsection{Autori}
	    \begin{itemize}
	        \item Dominik Boboš (\href{mailto:xbobos00@stud.fit.vutbr.cz}{xbobos00@stud.fit.vutbr.cz})
	        \item Dávid Oravec (\href{mailto:xorave05@stud.fit.vutbr.cz}{xorave05@stud.fit.vutbr.cz})
	    \end{itemize}
	    
	    \subsubsection{Literatúra}
	    Ako primárny zdroj pre našu štúdiu a zároveň \uv{odrazový mostík} sme použili \cite{maxbrenner2020} spolu s literárnymi zdrojmi na ktoré sa v štúdii odkazuje.
	    
	    Ostatné literárne zdroje sú uvedené v sekcii \hyperref[sec:bib]{Bibliografia}
	
	
	\newpage
    \bibliographystyle{czechiso}
    \renewcommand{\refname}{Bibliografia}
    \label{sec:bib}
    \bibliography{ims}

\end{document}
